% Kurzer Überblick: Kontext des Projekts (FRoST-Endeffektor), Motivation, Ziel
% (Positions- und Drehmomentregelung), und Aufbau des Berichts (Kapitelübersicht).

Die vorliegende Projektarbeit, die als Prüfungsleistung des Moduls ``Simulation und Regelung'' durchzuführen ist, lag eingangs der Thematik der Entwicklung eines Positions- und Drehmomentreglers für einen neu entwickelten FRoST-Endeffektor zugrunde. Aufgrund des Verwerfens des neuen Greifer-Konzeptes, wurde die Aufgabenstellung im Laufe der Bearbeitungszeit angepasst bzw. erweitert. Aufgrund der Absenz eines physischen Greifers, konnten Teilaspekte der Aufgabenstellung, die eine Vermessung der Mechanik voraussetzen, nicht bearbeitet werden. Zudem wurde die Entwicklung eines Prüfstandes, der die Funktionalität eines Greifers simuliert, als weiterer Bestandteil der Projektarbeit aufgenommen.\\

Das grundsätzlich (unveränderte) Ziel der Projektarbeit liegt darin, zwei Reglerkonzepte für einen Schrittmotor - dieser treibt die Finger an - zu entwickeln und praktisch umzusetzen. Es liegen lediglich vereinzelte Vorgaben hinsichtlich der verwendeten Hard- und Software vor. Durch den Positionsregler soll die Greifposition und durch den Drehmomentregler soll die Greifkraft sicher und präzise gesteuert werden können.\\

In den folgenden Kapiteln werden zuerst die Aufgabenstellung und die Anforderungen an das zu realisierende Regelsystem zusammengefasst (Kapitel \ref{Kap:Aufgabe}). Anschließend wird in Kapitel \ref{Kap:Material} die ausgewählte Hard- und Software vorgestellt. In Kapitel \ref{Kap:System} wird nachfolgend die Erarbeitung eines Entwurfs eines Prüfstandes als zu regelndes System und die daraus abgeleiteten Systemgrenzen und das daraus resultierende mathematische Modell beschrieben. In Kapitel \ref{Kap:Feedback} wird daraufhin das Auslesen des Motorfeedbacks mittels geeigneter Sensorik thematisiert. Infolgedessen werden die Entwürfe der beiden Regler in Kapitel \ref{Kap:Regler} erläutert. Die Implementierung der Firmware für das zentrale Steuergerät des Systems wird in Kapitel \ref{Kap:Firmware} beleuchtet. Um die entwickelten Regler auf ihre Funktionalität zu überprüfen, wird in Kapitel \ref{Kap:Testlauf} eingangs die Realisierung des Prüfstandes behandelt und anschließend die Einhaltung der vorgegebenen Anforderungen erörtert. Abschließend beinhaltet das Kapitel \ref{Kap:Fazit} ein Fazit und einen Ausblick.

